\documentclass[10pt]{article}
\usepackage[T1]{fontenc}
\usepackage[utf8]{inputenc}
\usepackage[french]{babel}
\usepackage{amsmath,amssymb,mathrsfs,mathpazo}
\usepackage{geometry}
\geometry{a4paper,margin=2cm}
\usepackage{graphicx,array,tabularx,tabularray,colortbl,multirow,multicol,booktabs,enumitem,pifont,ulem,fancybox,diagbox,cancel}
\usepackage{pstricks,pstricks-add,pst-plot,pst-tree,pst-node,pst-text,pst-func,pst-math,pst-bspline,pst-3dplot,pst-eucl,pst-solides3d}
\let\pscalculate\undefined
\usepackage{pst-bezier}
\usepackage{tikz,pgfplots,tkz-tab}
\usetikzlibrary{arrows,arrows.meta,patterns,calc,positioning,trees,decorations.markings,shapes.geometric,shapes.misc,shapes.arrows,intersections,decorations.pathmorphing,decorations.pathreplacing,backgrounds,shadows,babel}
\pgfplotsset{compat=1.18}
\usepgfplotslibrary{fillbetween}
\usepackage[np]{numprint}
\usepackage{listings,makecell,pas-tableur,twoopt}
\let\barre\undefined
\usepackage[pstricks]{tablvar}
\definecolor{darkgreen}{RGB}{0,100,0}
\providecommand{\red}{\color{red}}
\providecommand{\blue}{\color{blue}}
\providecommand{\green}{\color{green}}
\providecommand{\black}{\color{black}}
\providecommand{\white}{\color{white}}
\providecommand{\gray}{\color{gray}}
\providecommand{\cyan}{\color{cyan}}
\providecommand{\magenta}{\color{magenta}}
\providecommand{\yellow}{\color{yellow}}
\providecommand{\violet}{\color{violet}}
\providecommand{\orange}{\color{orange}}
\providecommand{\euro}{\texteuro}
\providecommand{\EUR}{\texteuro}
\providecommand{\ieme}{\textsuperscript{e}}
\providecommand{\degres}{\ensuremath{^\circ}}
\providecommand{\degree}{\ensuremath{^\circ}}
\providecommand{\vv}[1]{\overrightarrow{#1}}
\providecommand{\Coord}[1]{#1}
\providecommand{\decosix}{}
\providecommand{\starredbullet}{\ensuremath{\bullet}}
\pagestyle{empty}
\setlength{\parindent}{0pt}
\usepackage{adjustbox}\DeclareUnicodeCharacter{2236}{:}\DeclareUnicodeCharacter{2190}{\ensuremath{\leftarrow}}
\providecommand{\R}{\mathbb{R}}
\providecommand{\N}{\mathbb{N}}
\providecommand{\Z}{\mathbb{Z}}
\providecommand{\Q}{\mathbb{Q}}
\providecommand{\D}{\mathbb{D}}
\providecommand{\C}{\mathbb{C}}
\providecommand{\e}{{\mathrm{e}}}
\providecommand{\vect}[1]{\overrightarrow{#1}}
\providecommand{\vectt}[1]{\overrightarrow{\mathrm{#1}}}
\providecommand{\vvt}[1]{\overrightarrow{\mathrm{#1}}}
\providecommand{\np}[2][]{{#2~#1}}
\providecommand{\vv}[1]{{\overrightarrow{#1}}}
\renewcommand{\d}{{\mathrm{d}}}
\providecommand{\degres}{{^\circ}}
\providecommand{\degre}{{^\circ}}
\providecommand{\blue}{\color{blue}}
\providecommand{\red}{\color{red}}
\providecommand{\enskip}{\quad}
\providecommand{\llbracket}{\lbrack\!\lbrack}
\providecommand{\rrbracket}{\rbrack\!\rbrack}
\providecommand{\euro}{{\text{€}}}
\providecommand{\no}{{\text{n°}}}
\providecommand{\interval}[2]{{[#1;#2]}}
\providecommand{\intervFO}[2]{{[#1;#2[}}
\providecommand{\dotfill}{\ldots}

\providecommand{\ds}{\displaystyle}
\providecommand{\textding}[1]{#1}
\providecommand{\Oij}{\(\mathrm{O};\vec{\imath},\vec{\jmath}\)}
\providecommand{\Oijk}{\(\mathrm{O};\vec{\imath},\vec{\jmath},\vec{k}\)}
\providecommand{\Ouv}{\(\mathrm{O};\vec{u},\vec{v}\)}
\providecommand{\intervFF}[2]{\ensuremath{[#1;#2]}}
\renewcommand{\euro}{\ifmmode\text{\texteuro}\else\texteuro\fi}
\DeclareUnicodeCharacter{00B0}{\ifmmode^{\circ}\else\textdegree\fi}
\DeclareUnicodeCharacter{20AC}{\ifmmode\text{\texteuro}\else\texteuro\fi}
% Énoncés : Éducation nationale. Transcription et corrigés : APMEP, Denis Vergès et contributeurs. Découpage et classement : CoopMaths, Sébastien Lozano, Éric Elter et contributeurs. Adaptation de la présentation : Maths Explorer.
% CC BY-SA : https://coopmaths.fr/www/about/
% https://coopmaths.fr/alea/static/eam/2026/tex/eam_specifique_2026_06_etrangers_2_cor.tex
\begin{document}
\noindent\textit{APMEP / CoopMaths. CC BY-SA. Présentation adaptée par Maths Explorer.}\par\medskip
Soit la fonction $f$ définie et dérivable sur $[0\,;\,10]$ par :\quad $f(x)=-x^{3}+4{,}5\,x^{2}-6x+2$.
	
	\begin{enumerate}
		\item Calculer l'expression de la dérivée $f'$ en fonction de $x$.
		
		\begin{color}{blue}
		On utilise la dérivation d'une somme de fonctions dérivables.
		
		$f'(x)=(-1)\times 3x^{2} + 4{,}5\times 2x - 6  = -3x^2 + 9x - 6$
		
		\end{color}
		
		
		\item Vérifier que, pour $x$ appartenant à $[0\,;\,10]$, $f'(x)=(3x-6)(1-x)$.

		\begin{color}{blue}
		On développe l'expression :
		
		\(
		(3x-6)(1-x) = 3x\times 1 +3x\times (-x) + (-6)\times 1 (-6)\times (-x) 
		\)
		
		On réduit les produits :
		
		\(
		(3x-6)(1-x) = 3x -3x^2 - 6 + 6x 
		\)
		
		
		On réduit les sommes :

		\(
		(3x-6)(1-x) = -3x^2 + 9x - 6  
		\)
		
				
		\end{color}
		
		\item Étudier le signe de $f'(x)$ pour $x$ appartenant à $[0\,;\,10]$.
		
		\begin{color}{blue}
		Pour étudier le signe de $f'(x)$, on étudie le signe de chacun des facteurs de l'expression factorisée de $f'(x)=(3x-6)(1-x)$.
		
		$3x - 6 = 0 \iff 3x = 6 \iff x = 2$
		
		$3x - 6 > 0 \iff 3x > 6 \iff x > 2$
		
		$1 - x = 0 \iff 1 = x \iff x = 1$
		
		$1 - x > 0 \iff 1 > x \iff x < 1$
		
		On forme le tableau de signes :
		
		
		\begin{center}
				\begin{tikzpicture}
					\tkzTabInit[lgt=3.5,espcl=1.4]
					{$x$/0.7,$3x-6$/0.7,$1-x$/0.7,signe de $f'(x)$/1.4}
					{$0$,$1$,$2$,$10$}
				\tkzTabLine{,-, ,+, z,+,}
				\tkzTabLine{,+, z,-,,-,}
				\tkzTabLine{,-,z,+,z,-,}
			\end{tikzpicture}
		\end{center}
		
		 
				
		\end{color}
		
		\item En déduire les variations de $f$ sur $[0\,;\,10]$.
		
		Les valeurs des éventuels maximums et minimums ne sont pas demandées.
		
		\begin{color}{blue}
		\begin{center}
				\begin{tikzpicture}
					\tkzTabInit[lgt=3.5,espcl=1.4]
					{$x$/0.7,signe de $f'(x)$/1.4,variation de $f$/2.4}
					{$0$,$1$,$2$,$10$}
				\tkzTabLine{,-,z,+,z,-,}
				\tkzTabVar{+/,-/,+/,-/}
			\end{tikzpicture}
		\end{center}
		\end{color}
		
	\end{enumerate}

\end{document}
